% Generated draft: requires mathematical and source review.
\documentclass[11pt]{article}
\usepackage[margin=25mm]{geometry}
\usepackage{fontspec}
\setmainfont{Latin Modern Roman}
\usepackage{amsmath,amssymb,mathtools,unicode-math}
\setmathfont{Latin Modern Math}
\usepackage{enumitem,needspace,xurl,hyperref}
\setlength{\parindent}{0pt}
\setlength{\parskip}{5pt}
\setlength{\emergencystretch}{3em}
\newcommand{\N}{\mathbb{N}}
\newcommand{\Z}{\mathbb{Z}}
\newcommand{\Q}{\mathbb{Q}}
\newcommand{\R}{\mathbb{R}}
\newcommand{\C}{\mathbb{C}}
\newcommand{\F}{\mathbb{F}}
\newcommand{\A}{\mathbb{A}}
\newcommand{\PP}{\mathbb{P}}
\newcommand{\E}{\mathbb{E}}
\newcommand{\eps}{\varepsilon}
\newcommand{\dd}{\,\mathrm{d}}
\begin{document}
\section*{UnsolvedMath ID: 701097}
\textbf{Source problem code:} catalogue entry
\subsection{Definitions and mathematical statement}
For each integer \(n\geq 3\), let \(F_n=\langle x_1,\ldots,x_n\rangle\) be the free group with the specified free basis, and let \(\mathbb S_n\) be the group of permutations of \(\{1,\ldots,n\}\). Define the group of conjugating automorphisms by
\[
C_n=\left\{\varphi\in\operatorname{Aut}(F_n):\ \exists\pi\in\mathbb S_n\ \exists(t_1,\ldots,t_n)\in F_n^n\ \forall i\in\{1,\ldots,n\},\quad \varphi(x_i)=t_i^{-1}x_{\pi(i)}t_i\right\},
\]
with the operation inherited from \(\operatorname{Aut}(F_n)\). Here the permutation and the conjugating elements may depend on \(\varphi\).

A group \(G\) is called linear if there exist a field \(K\), a positive integer \(d\), and an injective group homomorphism \(G\hookrightarrow\operatorname{GL}_d(K)\).

The question is whether \(C_n\) is linear for every integer \(n\geq 3\); explicitly, whether
\[
\forall n\in\Z,\ n\geq 3,\quad \exists\text{ a field }K_n\ \exists d_n\in\Z_{\geq 1}\ \exists\text{ an injective homomorphism }\rho_n:C_n\longrightarrow\operatorname{GL}_{d_n}(K_n).
\]
The source does not prescribe a field or its characteristic. This formulation uses the convention above and permits both the field and the dimension to depend on \(n\).

For context, the supplied source identifies the subgroup of conjugating automorphisms fixing \(x_1x_2\cdots x_n\) with the braid group \(B_n\), and states that \(B_n\) is linear for every \(n\geq 2\), whereas \(\operatorname{Aut}(F_n)\) is not linear for \(n\geq 3\). These are supplied background assertions, not conclusions established by this draft.
\subsection{Short English statement}
For each rank \(n\geq 3\), consider all automorphisms of the free group that send each specified free generator to a conjugate of a generator, allowing a permutation of the generators. Must this group admit a faithful representation by finite square matrices over some field?
\subsection{Sources}
\begin{itemize}
\item Local UnsolvedMath registry, permanent ID \(701097\), exact source problem code \texttt{catalogue entry}, title \emph{Kourovka Notebook Problem 15.9}. The supplied registry labels the problem open; this status has not been independently verified.
\item V. G. Bardakov, proposer, \emph{The Kourovka Notebook (Unsolved Problems in Group Theory)}, Problem 15.9, 15th issue (2002). The supplied record also identifies the 21st issue edition (2026), \texttt{arXiv:1401.0300v48}. The registry describes the collection as edited by E. I. Khukhro and V. D. Mazurov. These bibliographic details and the source attribution are supplied data and have not been independently checked.
\end{itemize}
\end{document}
